\documentclass[UTF8,a4paper,12pt,oneside,openany,fontset=fandol]{ctexbook}
\usepackage{amsmath,amssymb,graphicx,geometry,fancyhdr,enumitem,array}
\usepackage{tikz}
\usetikzlibrary{arrows.meta,calc,patterns,decorations.markings,angles,quotes}
\tikzset{
  ray/.style={thin,postaction={decorate},decoration={markings,mark=at position .55 with {\arrow{Stealth[length=4pt,width=3pt]}}}},
  rayend/.style={thin,-{Stealth[length=4pt,width=3pt]}},
  dim/.style={very thin,{Stealth[length=3.5pt,width=2.5pt]}-{Stealth[length=3.5pt,width=2.5pt]}},
  axis/.style={thin},
  dashedaux/.style={thin,dashed}
}
\usepackage[unicode,hidelinks]{hyperref}

\geometry{left=22mm,right=22mm,top=20mm,bottom=20mm,headheight=16pt}
\pagestyle{fancy}
\fancyhf{}
\fancyhead[L]{\small 姚启钧光学教程做题本}
\fancyhead[R]{\small\nouppercase{\leftmark}}
\fancyfoot[C]{\thepage}
\renewcommand{\chaptermark}[1]{\markboth{第\thechapter 章\quad #1}{}}

\setlength{\parindent}{2em}
\setlength{\emergencystretch}{3em}
\linespread{1.12}
\ctexset{chapter={beforeskip=8pt,afterskip=16pt}}
\setlist[enumerate]{label=（\arabic*）,leftmargin=2.5em,itemsep=2pt,topsep=3pt}

\newcommand{\problem}[2]{%
  \par\phantomsection\pdfbookmark[1]{题 #1}{problem-#1}%
  \noindent\textbf{#1}\quad #2\par
}
\newcommand{\sproblem}[2]{%
  \par\phantomsection\pdfbookmark[1]{*题 #1}{problem-#1}%
  \noindent\textbf{\textsuperscript{*}#1}\quad #2\par
}

\newsavebox{\figbox}
\newlength{\maxfigheight}
\newcommand{\drawfig}[2]{%
  \begingroup
  \setlength{\maxfigheight}{42mm}%
  \sbox{\figbox}{\resizebox{#1\linewidth}{!}{\input{tikzfigs/fig_#2.tex}}}%
  \ifdim\dimexpr\ht\figbox+\dp\figbox\relax>\maxfigheight
    \resizebox{!}{\maxfigheight}{\usebox{\figbox}}%
  \else
    \usebox{\figbox}%
  \fi
  \endgroup
}
\def\figlabel#1_#2\relax{#1.#2}
\newcommand{\fig}[2][0.48]{%
  \par\addvspace{4pt}%
  {\centering\drawfig{#1}{#2}\par\vspace{2pt}{\footnotesize 题\figlabel#2\relax\ 图}\par}%
  \addvspace{2pt}%
}
\newcommand{\figcap}[3][0.48]{%
  \par\addvspace{4pt}%
  {\centering\drawfig{#1}{#2}\par\vspace{2pt}{\footnotesize #3}\par}%
  \addvspace{2pt}%
}
\newcommand{\problemfig}[4][0.34]{%
  \par\phantomsection\pdfbookmark[1]{题 #2}{problem-#2}%
  \noindent\begin{minipage}[t]{\dimexpr\linewidth-#1\linewidth-6mm\relax}%
    \vspace{0pt}\textbf{#2}\quad #3%
  \end{minipage}\hfill%
  \begin{minipage}[t]{#1\linewidth}%
    \vspace{0pt}\centering\drawfig{1}{#4}\par\vspace{2pt}{\footnotesize 题\figlabel#4\relax\ 图}%
  \end{minipage}\par
}
\newcommand{\sproblemfig}[4][0.34]{%
  \par\phantomsection\pdfbookmark[1]{*题 #2}{problem-#2}%
  \noindent\begin{minipage}[t]{\dimexpr\linewidth-#1\linewidth-6mm\relax}%
    \vspace{0pt}\textbf{\textsuperscript{*}#2}\quad #3%
  \end{minipage}\hfill%
  \begin{minipage}[t]{#1\linewidth}%
    \vspace{0pt}\centering\drawfig{1}{#4}\par\vspace{2pt}{\footnotesize 题\figlabel#4\relax\ 图}%
  \end{minipage}\par
}

\newcommand{\answerarea}{\par\vfill}

\begin{document}
\begin{titlepage}
  \centering
  \vspace*{.24\textheight}
  {\Huge\bfseries 姚启钧《光学教程》\\[0.6em]做题本\par}
  \vspace{1.8cm}
  {\Large 习题整理与作答空间\par}
  \vspace{0.8cm}
  {\large 第一至八章\quad 共 191 题\quad 每页两题\par}
  \vfill
  {\large 姓名：\underline{\hspace{4cm}}\quad 日期：\underline{\hspace{3cm}}\par}
\end{titlepage}

\frontmatter
\tableofcontents
\clearpage

\mainmatter
\chapter{光的干涉}

\problem{1.1}{波长为 $500\,\mathrm{nm}$ 的绿光投射在间距 $d$ 为 $0.022\,\mathrm{cm}$ 的双缝上，在距离 $r_0$ 为 $180\,\mathrm{cm}$ 处的光屏上形成干涉条纹，求两个亮条纹之间的距离。若改用波长为 $700\,\mathrm{nm}$ 的红光投射到此双缝上，两个亮纹之间的距离又为多少？算出这两种光第 2 级亮纹位置的距离。}
\answerarea

\problem{1.2}{在杨氏实验装置中，光源波长为 $640\,\mathrm{nm}$，两狭缝间距 $d$ 为 $0.4\,\mathrm{mm}$，光屏离狭缝的距离 $r_0$ 为 $50\,\mathrm{cm}$。试求：（1）光屏上第 1 级亮条纹和中央亮条纹之间的距离；（2）若 $P$ 点离中央亮条纹为 $0.1\,\mathrm{mm}$，问两束光在 $P$ 点的相位差是多少？（3）求 $P$ 点的光强度和中央点的强度之比。}
\answerarea
\clearpage

\problem{1.3}{把折射率为 $1.5$ 的玻璃片插入杨氏实验的一束光路中，光屏上原来第 5 级亮条纹所在的位置变为中央亮条纹，试求插入的玻璃片的厚度。已知光波长为 $600\,\mathrm{nm}$。}
\answerarea

\problem{1.4}{波长为 $500\,\mathrm{nm}$ 的单色平行光射在间距 $d$ 为 $0.2\,\mathrm{mm}$ 的双狭缝上，通过其中一个缝的能量为另一个的 $2$ 倍，在离狭缝 $50\,\mathrm{cm}$ 的光屏上形成干涉图样，求干涉条纹间距和条纹的可见度。}
\answerarea
\clearpage

\problem{1.5}{波长为 $700\,\mathrm{nm}$ 的光源与菲涅耳双镜的相交棱之间距离为 $20\,\mathrm{cm}$，棱到光屏间的距离 $L$ 为 $180\,\mathrm{cm}$，若所得干涉条纹中相邻亮条纹的间隔为 $1\,\mathrm{mm}$，求双镜平面之间的夹角 $\theta$。}
\answerarea

\problem{1.6}{在题 1.6 图所示的劳埃德镜实验中，光源 $S$ 到观察屏的距离为 $1.5\,\mathrm{m}$，到劳埃德镜面的垂直距离为 $2\,\mathrm{mm}$。劳埃德镜长 $40\,\mathrm{cm}$，置于光源和屏之间的中央。（1）若光波波长 $\lambda=500\,\mathrm{nm}$，问条纹间距是多少？（2）确定屏上可以看见条纹的区域大小，此区域内共有几条条纹？（提示：产生干涉的区域 $P_1P_2$ 可由图中的几何关系求得。）\fig[0.62]{1_6}}
\answerarea
\clearpage

\problem{1.7}{试求能产生红光（$\lambda=700\,\mathrm{nm}$）的二级反射亮干涉条纹的肥皂膜厚度，已知肥皂膜折射率为 $1.33$，且平行光与法向成 $30^\circ$ 角入射。}
\answerarea

\problem{1.8}{透镜表面通常镀一层如 $\mathrm{MgF_2}$（$n=1.38$）一类的透明物质薄膜，目的是利用干涉来降低玻璃表面的反射。为了使透镜在可见光谱的中心波长（$550\,\mathrm{nm}$）处产生极小的反射，则镀层最小厚度为多少？}
\answerarea
\clearpage

\problem{1.9}{在两块玻璃片之间一边放一条厚纸，另一边相互压紧。玻璃片 $l$ 长 $10\,\mathrm{cm}$，纸厚为 $0.05\,\mathrm{mm}$，从 $60^\circ$ 的反射角进行观察，问在玻璃片单位长度内看到的干涉条纹数目是多少？设单色光源波长为 $500\,\mathrm{nm}$。}
\answerarea

\problem{1.10}{在上题装置中，沿垂直于玻璃片表面的方向看去，看到相邻两条暗纹间距为 $1.4\,\mathrm{mm}$。已知玻璃片长 $17.9\,\mathrm{cm}$，纸厚 $0.036\,\mathrm{mm}$，求光波的波长。}
\answerarea
\clearpage

\problem{1.11}{波长为 $400\sim 760\,\mathrm{nm}$ 的可见光正射在一块厚度为 $1.2\times 10^{-6}\,\mathrm{m}$、折射率为 $1.5$ 的薄玻璃片上，试问从玻璃片反射的光中哪些波长的光最强。}
\answerarea

\problem{1.12}{迈克耳孙干涉仪的反射镜 $\mathrm{M}_2$ 移动 $0.25\,\mathrm{mm}$ 时，看到条纹移过的数目为 $909$ 个，设光为垂直入射，求所用光源的波长。}
\answerarea
\clearpage

\problem{1.13}{迈克耳孙干涉仪平面镜的面积为 $4\times 4\,\mathrm{cm}^2$，观察到该镜上有 $20$ 个条纹，当入射光的波长为 $589\,\mathrm{nm}$ 时，两镜面之间的夹角为多大？}
\answerarea

\problem{1.14}{调节一台迈克耳孙干涉仪，使其用波长为 $500\,\mathrm{nm}$ 的扩展光源照明时会出现同心圆环条纹，若要使圆环中心处相继出现 $1\,000$ 条圆环条纹，则必须将移动一臂多远的距离？若中心是亮的，试计算第一暗环的角半径。（提示：圆环是等倾干涉图样，计算第一暗环角半径时可利用 $\theta\approx\sin\theta$ 及 $\cos\theta\approx 1-\dfrac{\theta^2}{2}$ 的关系。）}
\answerarea
\clearpage

\problem{1.15}{用单色光观察牛顿环，测得 $j$ 级亮环的直径为 $3\,\mathrm{mm}$，在它的第 $j+5$ 级亮环的直径为 $4.6\,\mathrm{mm}$，所用平凸透镜的凸面曲率半径为 $1.03\,\mathrm{m}$，求此单色光的波长。}
\answerarea

\problem{1.16}{在反射光中观察某单色光所形成的牛顿环，其第 2 级亮环与第 3 级亮环间距为 $1\,\mathrm{mm}$，求第 19 和 20 级亮环之间的距离。}
\answerarea
\clearpage

\sproblemfig[0.32]{1.17}{牛顿环可由两个曲率半径很大的平凸透镜之间的空气层产生（题 1.17 图），平凸透镜 A 和 B 的曲率半径分别为 $R_{\mathrm{A}}$ 和 $R_{\mathrm{B}}$，在波长为 $600\,\mathrm{nm}$ 的单色光垂直照射下观察到反射牛顿环第 10 级暗环半径 $r_{\mathrm{AB}}=4\,\mathrm{mm}$。若另有曲率半径为 $R_{\mathrm{C}}$ 的平凸透镜 C（图中未画出），并且 B、C 组合和 A、C 组合产生的第 10 级暗环半径分别为 $r_{\mathrm{BC}}=4.5\,\mathrm{mm}$ 和 $r_{\mathrm{AC}}=5\,\mathrm{mm}$，试计算 $R_{\mathrm{A}}$、$R_{\mathrm{B}}$ 和 $R_{\mathrm{C}}$。}{1_17}
\answerarea

\sproblem{1.18}{菲涅耳双棱镜实验装置尺寸如下：缝到棱镜的距离为 $5\,\mathrm{cm}$，棱镜到屏的距离为 $95\,\mathrm{cm}$，棱镜角为 $\alpha=179^\circ 32'$，构成棱镜玻璃材料的折射率 $n'=1.5$，采用的是单色光。当厚度均匀的肥皂膜横过双棱镜的一半部分放置，该系统中心部分附近的条纹相对原先有 $0.8\,\mathrm{mm}$ 的位移。若肥皂膜的折射率为 $n=1.35$，试计算肥皂膜厚度的最小值为多少？}
\answerarea
\clearpage

\sproblemfig[0.18]{1.19}{将焦距为 $50\,\mathrm{cm}$ 的会聚透镜中央部分 C 切去（见题 1.19 图），余下的 A、B 两部分仍旧粘起来，C 的宽度为 $1\,\mathrm{cm}$。在对称轴线上距透镜 $25\,\mathrm{cm}$ 处置一点光源，发出波长为 $692\,\mathrm{nm}$ 的激光，在对称轴线上透镜的另一侧 $50\,\mathrm{cm}$ 处置一光屏，屏面垂直于轴线。试求：
\begin{enumerate}
  \item 干涉条纹的间距是多少？
  \item 光屏上呈现的干涉图样是怎样的？
\end{enumerate}}{1_19}
\answerarea

\sproblem{1.20}{将焦距为 $5\,\mathrm{cm}$ 的薄凸透镜 L 沿直线方向剖开（见题 1.20 图）分成两部分 A 和 B，并将 A 部分沿主轴右移至 $2.5\,\mathrm{cm}$ 处，这种类型的装置称为梅斯林对切透镜。若将波长为 $632.8\,\mathrm{nm}$ 的点光源 $P$ 置于主轴上离透镜 B 的距离为 $10\,\mathrm{cm}$ 处，试分析：
\begin{enumerate}
  \item 根据中学几何光学的知识分析成像情况如何？
  \item 若在 B 右边 $10.5\,\mathrm{cm}$ 处置一光屏，则在光屏上观察到的干涉图样如何？
\end{enumerate}\fig[0.45]{1_20}}
\answerarea
\clearpage

\sproblem{1.21}{如题 1.21 图所示，A 为平凸透镜，B 为平玻璃板，C 为金属柱，D 为框架，A、B 间有孔隙，图中绘出的是接触的情况，而 A 固结在框架的边缘上。温度变化时，C 发生伸缩，而假设 A、B、D 都不发生伸缩。以波长 $632.8\,\mathrm{nm}$ 的激光垂直照射，试问：
\begin{enumerate}
  \item 在反射光中观察时，看到牛顿环条纹移向中央，这表示金属柱 C 的长度在增加还是减小？
  \item 若观察到有 $10$ 个亮条纹移到中央而消失，试问 C 的长度变化了多少毫米？
\end{enumerate}\fig[0.44]{1_21}}
\answerarea

\sproblem{1.22}{研究性课题：请查阅互联网，进一步了解光的干涉应用的新进展，并写成评述性论文。}
\answerarea
\clearpage

\chapter{光的衍射}

\problem{2.1}{一平行单色光照射到一小圆孔上，将其波面分成半波带，求第 $k$ 个带的半径。若极点到观察点的距离 $r_0$ 为 $1\,\mathrm{m}$，单色光波长为 $450\,\mathrm{nm}$，求此时第一半波带的半径。}
\answerarea

\problem{2.2}{一平行单色光从左向右垂直射到一个有圆形小孔的屏上，设此孔可以像摄像机光圈那样改变大小。问：（1）小孔半径应满足什么条件时，才能使得此小孔右侧轴线上距小孔中心 $4\,\mathrm{m}$ 的 $P$ 点的光强分别得到极大值和极小值；（2）$P$ 点最亮时，小孔直径应为多大？设此光的波长为 $500\,\mathrm{nm}$。}
\answerarea
\clearpage

\problem{2.3}{波长为 $500\,\mathrm{nm}$ 的单色点光源离光阑 $1\,\mathrm{m}$，光阑上有一个内外半径分别为 $0.5\,\mathrm{mm}$ 和 $1\,\mathrm{mm}$ 的透光圆环，接收点 $P$ 离光阑 $1\,\mathrm{m}$，求 $P$ 点的光强 $I$ 与没有光阑时的光强度 $I_0$ 之比。}
\answerarea

\problem{2.4}{波长为 $632.8\,\mathrm{nm}$ 的 $\mathrm{He\text{-}Ne}$ 激光射向直径为 $2.76\,\mathrm{mm}$ 的圆孔，与孔相距 $1\,\mathrm{m}$ 处放一屏。试问：（1）屏上正对圆孔中心的 $P$ 点是亮点还是暗点？（2）要使 $P$ 点变成与（1）相反的情况，至少要把屏幕分别向前或向后移动多少？}
\answerarea
\clearpage

\problem{2.5}{一波带片由五个半波带组成。第一半波带为半径 $r_1$ 的不透明圆盘，第二半波带是半径 $r_1$ 至 $r_2$ 的透明圆环，第三半波带是 $r_2$ 至 $r_3$ 的不透明圆环，第四半波带是 $r_3$ 至 $r_4$ 的透明圆环，第五半波带是 $r_4$ 至无穷大的不透明区域。已知 $r_1:r_2:r_3:r_4=1:\sqrt{2}:\sqrt{3}:\sqrt{4}$，用波长 $500\,\mathrm{nm}$ 的平行单色光照明，最亮的像点在距波带片 $1\,\mathrm{m}$ 的轴上。试求：（1）$r_1$；（2）像点的光强与入射光强之比；（3）光强极大值出现在轴上哪些位置上。}
\answerarea

\problem{2.6}{波长为 $\lambda$ 的点光源经波带片成一个像点，该波带片有 $100$ 个透明奇数半波带（$1,3,5,\cdots,199$），另外 $100$ 个不透明偶数半波带。比较用波带片和换上同样焦距和口径的透镜时该像点的强度比 $I:I_0$。}
\answerarea
\clearpage

\problem{2.7}{一平行单色光的波长为 $480\,\mathrm{nm}$，垂直照射到宽度为 $0.4\,\mathrm{mm}$ 的狭缝上，会聚透镜的焦距为 $60\,\mathrm{cm}$。分别计算当缝的两边到 $P$ 点的相位差为 $\pi/2$ 和 $\pi/6$ 时，$P$ 点离焦点的距离。}
\answerarea

\problem{2.8}{白光形成的单缝衍射图样中，其中某一波长的第三个次最大值与波长为 $600\,\mathrm{nm}$ 的光波的第二个次最大值重合，求该光波的波长。}
\answerarea
\clearpage

\problem{2.9}{波长为 $546.1\,\mathrm{nm}$ 的平行单色光垂直地射在 $1\,\mathrm{mm}$ 宽的缝上，若将焦距为 $100\,\mathrm{cm}$ 的透镜紧贴于缝的后面，并使光聚焦到屏上，试问衍射图样的中央到（1）第一最小值；（2）第一最大值；（3）第三最小值的距离分别为多少？}
\answerarea

\problem{2.10}{钠光通过宽 $0.2\,\mathrm{mm}$ 的狭缝后，投射到与缝相距 $300\,\mathrm{cm}$ 的照相底片上，所得的第一最小值与第二最小值间的距离为 $0.885\,\mathrm{cm}$，试问钠光的波长为多少？若改用 X 射线（$\lambda=0.1\,\mathrm{nm}$）做此实验，问底片上这两个最小值之间的距离是多少？}
\answerarea
\clearpage

\problem{2.11}{以纵坐标表示强度，横坐标表示屏上的位置，粗略地画出三缝的夫琅禾费衍射（包括缝与缝之间的干涉）图样。设缝宽为 $b$，相邻缝间的距离为 $d$，$d=3b$。注意缺级问题。}
\answerarea

\problem{2.12}{一束平行白光垂直入射在每毫米 $50$ 条刻痕的全息光栅上，问第一级光谱的末端和第二光谱的始端的衍射角 $\theta$ 之差为多少？（设可见光中最短的紫光波长为 $400\,\mathrm{nm}$，最长的红光波长为 $760\,\mathrm{nm}$。）}
\answerarea
\clearpage

\problem{2.13}{用可见光（$760\sim 400\,\mathrm{nm}$）照射全息光栅时，一级光谱和二级光谱是否重叠？二级和三级怎样？若重叠，则重叠范围是多少？}
\answerarea

\problem{2.14}{用波长为 $589\,\mathrm{nm}$ 的平行单色光照射一衍射光栅，其光谱的中央最大值和第二十级主最大值之间的衍射角为 $15^\circ 10'$，求该光栅 $1\,\mathrm{cm}$ 内的缝数。}
\answerarea
\clearpage

\problem{2.15}{用每毫米内有 $400$ 条刻痕的平面透射光栅观察波长为 $589\,\mathrm{nm}$ 的钠光谱，试问：（1）光垂直入射时，最多能观察到几级光谱？（2）光以 $30^\circ$ 角入射时，最多能观察到几级光谱？}
\answerarea

\problem{2.16}{白光垂直照射到一个每毫米 $250$ 条刻痕的平面透射光栅上，试问在衍射角为 $30^\circ$ 处会出现哪些波长的光？其颜色如何？}
\answerarea
\clearpage

\problem{2.17}{一波长为 $624\,\mathrm{nm}$ 的平行单色光照射一光栅，已知该光栅的缝宽 $b$ 为 $0.012\,\mathrm{mm}$，不透明部分的宽度 $a$ 为 $0.029\,\mathrm{mm}$，缝数 $N$ 为 $10^3$ 条。试求：（1）单缝衍射图样的中央角宽度；（2）单缝衍射图样中央宽度内能看到多少级光谱？（3）谱线的半宽度为多少？}
\answerarea

\problem{2.18}{$\mathrm{NaCl}$ 的晶体结构是简单的立方点阵，其相对分子质量 $M=58.5$，密度 $\rho=2.17\,\mathrm{g/cm^3}$，（1）试证相邻两离子间的平均距离为
\[\sqrt[3]{\frac{M}{2N_{\mathrm{A}}\rho}}=0.281\,9\,\mathrm{nm}\]
式中 $N_{\mathrm{A}}=6.02\times 10^{23}\,\mathrm{mol^{-1}}$ 为阿伏伽德罗常量；（2）用 X 射线照射晶面时，第二级光谱的最大值在掠射角为 $1^\circ$ 的方向上出现，试计算该 X 射线的波长。}
\answerarea
\clearpage

\problem{2.19}{波长为 $0.001\,47\,\mathrm{nm}$ 的平行 X 射线射在晶体界面上，晶体原子层的间距为 $0.28\,\mathrm{nm}$，试问光线与界面成什么角度时，能观察到二级光谱。}
\answerarea

\sproblemfig[0.34]{2.20}{如题 2.20 图所示有三条彼此平行的狭缝，宽度均为 $b$，缝距分别为 $d$ 和 $2d$，试用振幅矢量叠加法证明正入射时，夫琅禾费衍射强度公式为
\[I_\theta = I_0\frac{\sin^2 u}{u^2}\bigl[3+2(\cos 2v+\cos 4v+\cos 6v)\bigr]\]
式中 $u=\dfrac{\pi b\sin\theta}{\lambda}$，$v=\dfrac{\pi d\sin\theta}{\lambda}$。}{2_20}
\answerarea
\clearpage

\sproblemfig[0.24]{2.21}{一宽度为 $2\,\mathrm{cm}$ 的平面透射光栅上刻有 $12\,000$ 条刻痕。如题 2.21 图所示，以波长 $\lambda=500\,\mathrm{nm}$ 的单色光垂直投射，将折射率为 $1.5$ 的劈状玻璃片置于光栅前方，玻璃片的厚度从光栅的一端到另一端由 $1\,\mathrm{mm}$ 均匀变薄到 $0.5\,\mathrm{mm}$，试问第一级主最大方向改变了多少？}{2_21}
\answerarea

\sproblem{2.22}{一平行单色光投射于平面透射光栅上，其方向与光栅的法线成 $\theta_0$ 角，在和法线成 $11^\circ$ 和 $53^\circ$ 角的方向上出现第一级谱线，且位于法线的两侧。
\begin{enumerate}
  \item 试求入射角 $\theta_0$；
  \item 试问为什么在法线两侧能观察到一级谱线，而在法线同侧则能观察到二级谱线？
\end{enumerate}}
\answerarea
\clearpage

\sproblem{2.23}{一波长为 $600\,\mathrm{nm}$ 的平行单色光正入射到一平面透射光栅上，有两个相邻主最大分别出现在 $\sin\theta_1=0.2$ 和 $\sin\theta=0.3$ 处，第四级为缺级。
\begin{enumerate}
  \item 试求光栅常量；
  \item 试求光栅的缝可能的最小宽度；
  \item 在确定了光栅常量与缝宽之后，试列出在光屏上实际呈现的全部级数。
\end{enumerate}}
\answerarea

\sproblem{2.24}{研究性课题：设计实验测定 CD 或 DVD 光盘的光栅常量。}
\answerarea
\clearpage

\chapter{几何光学的基本原理}

\problem{3.1}{证明反射定律符合费马原理。}
\answerarea

\sproblem{3.2}{根据费马原理可以导出在近轴光线条件下，从物点发出并会聚到像点的所有光线的光程都相等，由此导出薄透镜的物像公式。}
\answerarea
\clearpage

\problem{3.3}{眼睛 $\mathrm{E}$ 和物体 $PQ$ 之间有一块折射率为 $1.5$ 的玻璃平板（见题 3.3 图），平板的厚度 $d$ 为 $30\,\mathrm{cm}$。求物体 $PQ$ 的像 $P'Q'$ 与物体 $PQ$ 之间的距离 $d_2$ 为多少？\fig[0.52]{3_3}}
\answerarea

\problem{3.4}{玻璃棱镜的折射棱角 $A$ 为 $60^\circ$，对某一波长的光其折射率 $n$ 为 $1.6$。计算：（1）最小偏向角；（2）此时的入射角；（3）能使光线从 $A$ 角两侧透过棱镜的最小入射角。}
\answerarea
\clearpage

\problem{3.5}{题 3.5 图表示一种恒偏向棱镜，它相当于一个 $30^\circ\text{--}60^\circ\text{--}90^\circ$ 棱镜与一个 $45^\circ\text{--}45^\circ\text{--}90^\circ$ 棱镜按图示方式组合在一起。白光沿 $i$ 方向入射，旋转棱镜改变 $\theta_1$，从而使任意一种波长的光可以依次循着图示的路径传播，出射光线为 $r$。求证：如果 $\sin\theta_1=\dfrac{n}{2}$，则 $\theta_2=\theta_1$，且光束 $i$ 与 $r$ 相互垂直（这就是恒偏向棱镜名字的由来）。\fig[0.48]{3_5}}
\answerarea

\problem{3.6}{高 $5\,\mathrm{cm}$ 的物体距凹面镜顶点 $12\,\mathrm{cm}$，凹面镜的焦距是 $10\,\mathrm{cm}$，求像的位置及高度，并作光路图。}
\answerarea
\clearpage

\problem{3.7}{一个 $5\,\mathrm{cm}$ 高的物体放在球面镜前 $10\,\mathrm{cm}$ 处成 $1\,\mathrm{cm}$ 高的虚像。求：（1）此镜的曲率半径；（2）此镜是凸面镜还是凹面镜？}
\answerarea

\problem{3.8}{某观察者通过一块薄玻璃板去看在凸面镜中自己的像。他移动着玻璃板，使得在玻璃板中与在凸面镜中所看到的他眼睛的像重合在一起。若凸面镜的焦距为 $10\,\mathrm{cm}$，眼睛距凸面镜顶点的距离为 $40\,\mathrm{cm}$，问玻璃板与观察者眼睛之间的距离为多少？}
\answerarea
\clearpage

\problem{3.9}{物体位于凹面镜轴线上焦点之外，在焦点与凹镜面之间放一个与轴线垂直的两表面互相平行的玻璃板，其厚度为 $d$，折射率为 $n$。试证明：放入该玻璃板后使像移动的距离与把凹面镜向物体移动 $d(n-1)/n$ 的一段距离的效果相同。}
\answerarea

\problem{3.10}{欲使由无穷远发出的近轴光线通过透明球体并成像在右半球面的顶点处，问这透明球体的折射率应为多少？}
\answerarea
\clearpage

\problem{3.11}{有一折射率为 $1.5$、半径为 $4\,\mathrm{cm}$ 的玻璃球，物体在距球表面 $6\,\mathrm{cm}$ 处，求：（1）物所成的像到球心之间的距离；（2）像的横向放大率。}
\answerarea

\problem{3.12}{一个折射率为 $1.53$、直径为 $20\,\mathrm{cm}$ 的玻璃球内有两个小气泡，看上去一个恰好在球心，另一个从离观察者最近的方向看去，好像在表面与球心连线的中点。求两气泡的实际位置。}
\answerarea
\clearpage

\problem{3.13}{直径为 $1\,\mathrm{m}$ 的球形鱼缸的中心处有一条小鱼，若玻璃缸壁的影响可忽略不计，求缸外观察者所看到的小鱼的表观位置和横向放大率。}
\answerarea

\problem{3.14}{玻璃棒一端成半球形，其曲率半径为 $2\,\mathrm{cm}$。将它水平地浸入折射率为 $1.33$ 的水中，沿着棒的轴线离球面顶点 $8\,\mathrm{cm}$ 处的水中有一物体，利用计算和作图法求像的位置及横向放大率，并作光路图。}
\answerarea
\clearpage

\problem{3.15}{有两块玻璃薄透镜的两表面均各为凸球面及凹球面，其曲率半径为 $10\,\mathrm{cm}$。一物点在主轴上距镜 $20\,\mathrm{cm}$ 处，若物和镜均浸在水中，分别用作图法和计算法求像点的位置。设玻璃的折射率为 $1.5$，水的折射率为 $1.33$。}
\answerarea

\problem{3.16}{一凸透镜在空气中的焦距为 $40\,\mathrm{cm}$，在水中时焦距为 $136.8\,\mathrm{cm}$，问此透镜的折射率为多少（水的折射率为 $1.33$）？若将此透镜置于 $\mathrm{CS_2}$ 中（$\mathrm{CS_2}$ 的折射率为 $1.62$），其焦距又为多少？}
\answerarea
\clearpage

\problem{3.17}{两片极薄的表面玻璃，曲率半径分别为 $20\,\mathrm{cm}$ 和 $25\,\mathrm{cm}$。将两片的边缘粘起来，形成内含空气的双凸透镜，把它置于水中，求其焦距为多少？}
\answerarea

\problem{3.18}{会聚透镜和发散透镜的焦距都是 $10\,\mathrm{cm}$，求：（1）与主轴成 $30^\circ$ 的一束平行光入射到每个透镜上，像点在何处？（2）在每个透镜左方的焦平面上离主轴 $1\,\mathrm{cm}$ 处各置一发光点，成像在何处？作出光路图。}
\answerarea
\clearpage

\problem{3.19}{题 3.19 图（a）、（b）所示的 $MM'$ 分别为一薄透镜的主光轴，$S$ 为光源，$S'$ 为像。用作图法求透镜中心和透镜焦点的位置。\fig[0.64]{3_19}}
\answerarea

\sproblem{3.20}{比累对切透镜是把一块凸透镜沿直径方向剖成两半组成，两半块透镜垂直光轴拉开一点距离，用挡光的光阑 $\mathrm{K}$ 挡住其间的空隙（见题 3.20 图），这时可在屏上观察到干涉条纹。已知点光源 $P$ 与透镜相距 $300\,\mathrm{cm}$，透镜的焦距 $f'=50\,\mathrm{cm}$，两半透镜拉开的距离 $t=1\,\mathrm{mm}$，光屏与透镜相距 $l=450\,\mathrm{cm}$。用波长为 $632.8\,\mathrm{nm}$ 的氦氖激光作为光源，求干涉条纹的间距。\fig[0.52]{3_20}}
\answerarea
\clearpage

\sproblemfig[0.38]{3.21}{把焦距为 $10\,\mathrm{cm}$ 的会聚透镜的中央部分 C 切去，C 的宽度为 $1\,\mathrm{cm}$，把余下的两部分粘起来（见题 3.21 图）。如在其对称轴上距透镜 $5\,\mathrm{cm}$ 处置一点光源，试求像的位置。}{3_21}
\answerarea

\sproblemfig[0.18]{3.22}{一折射率为 $1.5$ 的薄透镜，其凸面的曲率半径为 $5\,\mathrm{cm}$，凹面的曲率半径为 $15\,\mathrm{cm}$，且镀上银（见题 3.22 图）。试证明：当光从凸表面入射时，该透镜的作用相当于一个平面镜。（提示：物经过凸面折射、凹面反射和凸面再次折射后，$s'=-s$，$\beta=1$。）}{3_22}
\answerarea
\clearpage

\problem{3.23}{题 3.23 图所示的是一个等边直角棱镜和两个透镜所组成的光学系统，棱镜折射率为 $1.5$，凸透镜的焦距为 $20\,\mathrm{cm}$，凹透镜的焦距为 $10\,\mathrm{cm}$，两透镜间距为 $5\,\mathrm{cm}$，凸透镜距棱镜边的距离为 $10\,\mathrm{cm}$。求图中长度为 $1\,\mathrm{cm}$ 的物体所成像的位置和大小。（提示：物经棱镜成像在透镜轴上，相当于经过一块厚 $6\,\mathrm{cm}$ 的平板玻璃，可利用例 3.1 的结果求棱镜所成像的位置。）\fig[0.50]{3_23}}
\answerarea

\problem{3.24}{显微镜由焦距为 $1\,\mathrm{cm}$ 的物镜和焦距为 $3\,\mathrm{cm}$ 的目镜组成，物镜与目镜之间的距离为 $20\,\mathrm{cm}$，问物体放在何处时才能使最后的像成在距离眼睛 $25\,\mathrm{cm}$ 处？作出光路图。}
\answerarea
\clearpage

\problemfig[0.40]{3.25}{题 3.25 图中 $\mathrm{L}$ 为薄透镜，水平横线 $MM'$ 为主轴，$ABC$ 为已知的一条穿过这个透镜的光线的路径，用作图法求出任一条光线 $DE$ 穿过透镜后的路径。}{3_25}
\answerarea

\sproblemfig[0.40]{3.26}{题 3.26 图中 $MM'$ 是一厚透镜的主轴，$H$、$H'$ 是透镜的主平面，$S_1$ 是点光源，$S'_1$ 是点光源的像。试用作图法求任意一物点 $S_2$ 的像 $S'_2$ 的位置。}{3_26}
\answerarea
\clearpage

\problem{3.27}{双凸薄透镜的折射率为 $1.5$，$|r_1|=10\,\mathrm{cm}$，$|r_2|=15\,\mathrm{cm}$，$r_2$ 的一面镀银，物点 $P$ 在透镜前主轴上 $20\,\mathrm{cm}$ 处，求最后像的位置并作出光路图。}
\answerarea

\problem{3.28}{实物与光屏间的距离为 $l$，在中间某一位置放一个凸透镜，可使实物的像清晰地投于屏上。将透镜移过距离 $d$ 之后，屏上又出现一个清晰的像。（1）试计算两个像的大小之比；（2）证明透镜的焦距为 $(l^2-d^2)/4l$；（3）证明 $l$ 不能小于透镜焦距的 $4$ 倍。}
\answerarea
\clearpage

\problem{3.29}{一厚透镜的焦距 $f'$ 为 $60\,\mathrm{mm}$，其两焦点间的距离为 $125\,\mathrm{mm}$，若（1）物点置于光轴上物方焦点左方 $20\,\mathrm{mm}$ 处；（2）物点置于光轴上物方焦点右方 $20\,\mathrm{mm}$ 处；（3）虚物落在光轴上像方主点右方 $20\,\mathrm{mm}$ 处，问在这三种情况下像的位置各在何处？像的性质各如何？并作光路图。}
\answerarea

\sproblem{3.30}{一个会聚薄透镜和一个发散薄透镜互相接触而成一复合光具组，当物距为 $-80\,\mathrm{cm}$ 时，实像距镜 $60\,\mathrm{cm}$，若会聚透镜的焦距为 $10\,\mathrm{cm}$，问发散透镜的焦距为多少？}
\answerarea
\clearpage

\sproblem{3.31}{双凸厚透镜两个球面表面的曲率半径分别为 $100\,\mathrm{mm}$ 和 $200\,\mathrm{mm}$，沿轴厚度为 $10\,\mathrm{mm}$，玻璃的折射率为 $1.5$，试求其焦点、主点的位置，并绘图表示之。}
\answerarea

\sproblem{3.32}{一条光线射到折射率为 $n$ 的一球形水滴（见题 3.32 图），求：（1）若后表面的入射角为 $\alpha$，问这条光线将被全反射还是部分反射？（2）偏转角 $\delta$；（3）产生最小偏转角的入射角 $\varphi$。\fig[0.52]{3_32}}
\answerarea
\clearpage

\sproblem{3.33}{将灯丝置于空心玻璃球的中心，玻璃球的内外直径分别为 $8\,\mathrm{cm}$ 和 $9\,\mathrm{cm}$。（1）求从球外观察到的灯丝像的位置（设玻璃折射率 $n=1.5$）；（2）玻璃温度计管子的内外直径分别为 $1\,\mathrm{mm}$ 和 $3\,\mathrm{mm}$，求从侧面观察到的直径的数值；（3）同一温度计竖直悬挂于直径 $100\,\mathrm{mm}$ 的盛水玻璃烧杯的正中，从较远处通过烧杯壁观察时，温度计的内、外直径分别为多少？}
\answerarea

\sproblem{3.34}{如题 3.34 图所示为梅斯林分波面干涉实验装置。其中 $O_1$、$O_2$ 分别为两块半透镜 $\mathrm{L_1}$ 和 $\mathrm{L_2}$ 的光心，$S$、$O_1$、$O_2$、$S_1$、$S_2$ 共轴，且 $S_1S_2=l$。（1）试证来自 $\mathrm{L_1}$ 和 $\mathrm{L_2}$ 两端的光束到达 $P$ 点的光程差 $\delta=l-(S_1P+S_2P)$；（2）定性讨论与轴线垂直的光屏上接收到的干涉图样的特点。\fig[0.56]{3_34}}
\answerarea
\clearpage

\sproblem{3.35}{把杂质扩散到玻璃中可以增大玻璃的折射率，这就有可能造出一个厚度均匀的透镜。已知圆板半径为 $R$，厚度为 $d$（如题 3.35 图所示），求沿半径变化的折射率 $n(r)$，它会使从 $A$ 点发出的光线传播到 $B$ 点。假定这是个薄透镜，$d\ll a$，$d\ll b$。\fig[0.52]{3_35}}
\answerarea

\sproblem{3.36}{一弯凸薄透镜两个表面的半径 $r_1$ 和 $r_2$ 分别为 $-20\,\mathrm{cm}$ 和 $-15\,\mathrm{cm}$，折射率为 $1.5$，在 $r_2$ 的凸面镀银。在距 $r_1$ 球面左侧 $40\,\mathrm{cm}$ 处的主轴上置一高为 $1\,\mathrm{cm}$ 的物，试求最后成像的位置和像的性质。}
\answerarea
\clearpage

\sproblem{3.37}{一折射率为 $n$，曲率半径为 $R_1$ 和 $R_2$ 的薄凸透镜放在折射率分别为 $n_1$ 和 $n_2$ 的两种介质之间，$s_1$ 和 $s_2$ 分别为物距和像距，$f_1$ 和 $f_2$ 是物方和像方焦距。证明：
\[\frac{f_1}{s_1}+\frac{f_2}{s_2}=1.\]}
\answerarea

\sproblem{3.38}{研究性课题：请查阅互联网，了解超薄纳米光学器件的新进展，并写成一篇小论文。}
\answerarea
\clearpage

\chapter{光学仪器的基本原理}

\problem{4.1}{眼睛的构造可简化为一折射球面模型，其曲率半径为 $5.55\,\mathrm{mm}$，内部为折射率等于 $4/3$ 的液体，外部是空气，其折射率近似地等于 $1$。试计算眼球的两个焦距。用肉眼观察月球时，月球对眼的张角为 $1^\circ$，问视网膜上月球的像有多大？}
\answerarea

\problem{4.2}{把人眼的晶状体简化成距视网膜 $2\,\mathrm{cm}$ 的一个凸透镜，有人能看清距离在 $100\,\mathrm{cm}$ 到 $300\,\mathrm{cm}$ 间的物体。试问：（1）此人看清远点和近点时，眼睛透镜的焦距是多少？（2）为看清 $25\,\mathrm{cm}$ 远的物体，需配戴怎样的眼镜？}
\answerarea
\clearpage

\problem{4.3}{一照相机对准远物时，底片距物镜 $18\,\mathrm{cm}$，当镜头拉至最大长度时，底片与物镜相距 $20\,\mathrm{cm}$，求目的物在镜头前的最近距离？}
\answerarea

\problem{4.4}{两星所成的视角为 $4'$，用望远镜物镜照相，所得两像点相距 $1\,\mathrm{mm}$，问望远镜物镜的焦距是多少？}
\answerarea
\clearpage

\problem{4.5}{一显微镜具有三个物镜和两个目镜，三个物镜的焦距分别为 $16\,\mathrm{mm}$、$4\,\mathrm{mm}$ 和 $1.9\,\mathrm{mm}$，两个目镜的标示分别为 $5\times$ 和 $10\times$。设三个物镜造成的像都能落在像距为 $160\,\mathrm{mm}$ 处，问这显微镜的最大和最小的放大本领各为多少？}
\answerarea

\problem{4.6}{一显微镜物镜焦距为 $0.5\,\mathrm{cm}$，目镜焦距为 $2\,\mathrm{cm}$，两镜间距为 $22\,\mathrm{cm}$。观察者看到的像在无穷远处，试求物体到物镜的距离和显微镜的放大本领。}
\answerarea
\clearpage

\problem{4.7}{试证明望远镜的物镜为有效光阑（假定物镜和目镜的孔径相差不太悬殊）。}
\answerarea

\problem{4.8}{已知望远镜物镜的边缘即为有效光阑，试计算并作图求入射光瞳和出射光瞳的位置。}
\answerarea
\clearpage

\problem{4.9}{组成题 4.9 图的简单望远镜中各薄透镜的参数为 $\mathrm{L_1}$：$f_1=10\,\mathrm{cm}$，$D_1=4\,\mathrm{cm}$；$\mathrm{L_2}$：$f_2=2\,\mathrm{cm}$，$D_2=1.2\,\mathrm{cm}$；$\mathrm{L_3}$：$f_3=2\,\mathrm{cm}$，$D_3=1.2\,\mathrm{cm}$。计算该系统出射光瞳的位置和大小。\fig[0.54]{4_9}}
\answerarea

\problem{4.10}{有一光阑直径为 $5\,\mathrm{cm}$，放置在薄透镜后 $3\,\mathrm{cm}$ 处。透镜的焦距为 $5\,\mathrm{cm}$，孔径为 $6\,\mathrm{cm}$。现有一高为 $3\,\mathrm{cm}$ 的物 $PQ$ 置于透镜前 $12\,\mathrm{cm}$ 处。试求：（1）计算对主轴上 $P$ 点的入射光瞳和出射光瞳的大小和位置；（2）找到像的位置；（3）作光路图。}
\answerarea
\clearpage

\problem{4.11}{题 4.11 图中的 $H$、$H'$ 为光具组的主点，$F$、$F'$ 为焦点，$E$ 为对于物点 $P$ 的入射光瞳，$EO$ 为其半径。已知 $EO=2\,\mathrm{cm}$，$HP=20\,\mathrm{cm}$，$HF=15\,\mathrm{cm}$，$HO=5\,\mathrm{cm}$，$H'F'=15\,\mathrm{cm}$，$HH'=5\,\mathrm{cm}$，物长 $PQ=0.5\,\mathrm{cm}$。试作光路图并计算：（1）像的位置；（2）像长；（3）入射孔径角；（4）$P$ 点的出射光瞳半径和出射孔径角。\fig[0.52]{4_11}}
\answerarea

\problem{4.12}{一可认为是点光源的节能灯悬在圆桌中央的上空，桌的半径为 $R$，为了使桌的边缘能得到最大的照度，灯应悬在离桌面中心多高处？}
\answerarea
\clearpage

\problem{4.13}{一焦距为 $20\,\mathrm{cm}$ 的薄透镜，放在发光强度为 $15\,\mathrm{cd}$ 的点光源之前 $30\,\mathrm{cm}$ 处，在透镜后面 $80\,\mathrm{cm}$ 处放一屏，在屏上得到明亮的圆斑。求不计透镜中光的吸收时，圆斑的中心照度。}
\answerarea

\problem{4.14}{一长为 $5\,\mathrm{mm}$ 的线状物体放在一数码相机镜头前 $50\,\mathrm{cm}$ 处，在 CCD 上形成的像长为 $1\,\mathrm{mm}$。若 CCD 后移 $1\,\mathrm{cm}$，则像的弥散斑宽度为 $1\,\mathrm{mm}$。试求数码相机镜头的 $F$ 数。}
\answerarea
\clearpage

\problem{4.15}{某种玻璃在靠近钠光的黄色双谱线（其波长分别为 $589\,\mathrm{nm}$ 和 $589.6\,\mathrm{nm}$）附近的色散率 $\mathrm{d}n/\mathrm{d}\lambda$ 为 $-360\,\mathrm{cm^{-1}}$，求由此种玻璃制成的能分辨钠光双谱线的三棱镜，底边宽度应不小于多少？}
\answerarea

\problem{4.16}{设计一块平面透射光栅，要求满足：（1）使波长 $600\,\mathrm{nm}$ 的第二级谱线的衍射角小于 $30^\circ$，并能分辨 $0.02\,\mathrm{nm}$ 的波长差；（2）色散尽可能大；（3）第三级谱线缺级。试求：（1）光栅的缝宽、缝数、光栅常量和总宽度；（2）用这块光栅总共能看到 $600\,\mathrm{nm}$ 的几条谱线？}
\answerarea
\clearpage

\problem{4.17}{若要求显微镜能分辨相距 $0.000\,375\,\mathrm{mm}$ 的两点，用波长为 $550\,\mathrm{nm}$ 的可见光照明。试求：（1）此显微镜物镜的数值孔径 $R_{\mathrm{N.A.}}$；（2）若要求此两点放大后的视角为 $2'$，则显微镜的放大本领是多少？}
\answerarea

\problem{4.18}{夜间自远处驶来轿车的两前灯相距 $1.5\,\mathrm{m}$。如将眼睛的瞳孔看成产生衍射的圆孔，试估计视力正常的人在多远处才能分辨出光源是两个灯。设眼睛瞳孔的直径为 $3\,\mathrm{mm}$，设光源发出的光的波长 $\lambda$ 为 $550\,\mathrm{nm}$。}
\answerarea
\clearpage

\problem{4.19}{用孔径分别为 $20\,\mathrm{cm}$ 和 $160\,\mathrm{cm}$ 的两种望远镜能否分辨清月球上直径为 $500\,\mathrm{m}$ 的环形山？已知月球与地面的距离为地球半径的 $60$ 倍，而地球半径约为 $6\,370\,\mathrm{km}$。设光源发出的光的波长 $\lambda$ 为 $550\,\mathrm{nm}$。}
\answerarea

\problem{4.20}{电子显微镜的孔径角 $2u=8^\circ$，电子束的波长为 $0.1\,\mathrm{nm}$，试求：（1）它的最小分辨距离；（2）若人眼能分辨在明视距离处相距 $6.7\times 10^{-2}\,\mathrm{mm}$ 的两点，则该显微镜的放大倍数是多少？}
\answerarea
\clearpage

\sproblem{4.21}{平行光垂直投射于宽度为 $4\,\mathrm{cm}$ 的理想平面透射光栅上，已知在衍射角为 $60^\circ$ 的方向上的角色散率为 $0.5\times 10^{-2}\,\mathrm{rad/nm}$，试求光栅在该方向上最大处的分辨本领 $P$？}
\answerarea

\sproblem{4.22}{$\mathrm{He\text{-}Ne}$ 激光器发出波长为 $632.8\,\mathrm{nm}$、截面直径为 $2\,\mathrm{mm}$ 的激光束投向月球。已知月球和地面的距离为 $3.76\times 10^5\,\mathrm{km}$。试问：（1）在月球上得到的光斑的直径有多大？（2）如果这个激光经扩束器扩束后截面直径分别为 $2\,\mathrm{m}$ 和 $5\,\mathrm{m}$，再发向月球，试问在月球表面上的光斑直径各为多大？}
\answerarea
\clearpage

\sproblem{4.23}{宇宙飞船中的一架照相机在离地面 $200\,\mathrm{km}$ 的太空拍摄地面上的物体，如果要求它能分辨地面上相距 $1\,\mathrm{m}$ 的两点，照相机镜头的口径至少要多大？设光的波长为 $550\,\mathrm{nm}$。}
\answerarea

\sproblem{4.24}{已知月地距离约为 $3.8\times 10^5\,\mathrm{km}$，用上海天文台的口径为 $1.56\,\mathrm{m}$ 天体测量望远镜能分辨月球表面上两点的最小距离为多少？设光的波长为 $555\,\mathrm{nm}$。}
\answerarea
\clearpage

\sproblem{4.25}{为了分辨第二级钠光谱的双线（$589\,\mathrm{nm}$、$589.6\,\mathrm{nm}$），长度为 $15\,\mathrm{cm}$ 的平面光栅的光栅常量应为多少？}
\answerarea

\sproblem{4.26}{研究性课题：以 CD 或 DVD 光盘作为反射光栅，设计一光谱仪。}
\answerarea
\clearpage

\sproblem{4.27}{试绘出房门中设置的“猫眼”的结构及其光路图。}
\answerarea
\clearpage

\chapter{光的偏振}

\problem{5.1}{试确定下面两列光波
\begin{align*}
  \boldsymbol{E}_1 &= A_0\left[\boldsymbol{e}_x\cos(\omega t-kz)+\boldsymbol{e}_y\cos\left(\omega t-kz-\frac{\pi}{2}\right)\right] \\
  \boldsymbol{E}_2 &= A_0\left[\boldsymbol{e}_x\sin(\omega t-kz)+\boldsymbol{e}_y\sin\left(\omega t-kz-\frac{\pi}{2}\right)\right]
\end{align*}
的偏振态。}
\answerarea

\problem{5.2}{为了比较两个被自然光照射的表面的亮度，对其中一个表面直接进行观察，另一个表面通过两块偏振片来观察，两偏振片的透振方向的夹角为 $60^\circ$。若观察到两表面的亮度相同，则两表面实际的亮度比是多少？已知光通过每一块偏振片后损失入射光能量的 $10\%$。}
\answerarea
\clearpage

\problem{5.3}{两个尼科耳棱镜 $\mathrm{N_1}$ 和 $\mathrm{N_2}$ 主截面的夹角为 $60^\circ$，在它们之间放置另一个尼科耳棱镜 $\mathrm{N_3}$，让平行的自然光通过这个系统。假设各尼科耳棱镜对非常光均无吸收，试问 $\mathrm{N_3}$ 和 $\mathrm{N_1}$ 的主截面的夹角为何值时，通过系统的光强最大？设入射光强为 $I_0$，求此时所能通过的最大光强。}
\answerarea

\problem{5.4}{在两个正交的理想偏振片之间有一个偏振片以匀角速度 $\omega$ 绕光的传播方向旋转（见题 5.4 图），若入射的自然光强为 $I_0$，试证明透射光强为
\[I=\frac{I_0}{16}(1-\cos 4\omega t)\]\fig[0.52]{5_4}}
\answerarea
\clearpage

\problem{5.5}{线偏振光入射到折射率为 $1.732$ 的玻璃片上，入射角是 $60^\circ$，入射光的电矢量与入射面成 $30^\circ$ 角，求由分界面上反射的光强占入射光强的百分比。}
\answerarea

\problemfig[0.34]{5.6}{如题 5.6 图所示，自然光从空气入射到水面上，折射光再投向在水中倾斜放置的玻璃片上。若使从水面与玻璃片上的反射光均为线偏振光，求玻璃片和水面的夹角 $\alpha$（$n_{\text{玻}}=1.50$）。}{5_6}
\answerarea
\clearpage

\problem{5.7}{一束理想的线偏振光垂直入射到一方解石晶体上，它的振动面和主截面成 $30^\circ$ 角，两束折射光通过在方解石后面的一个尼科耳棱镜，其主截面与入射光的振动方向成 $50^\circ$ 角。计算两束透射光的相对强度。}
\answerarea

\problem{5.8}{线偏振光垂直入射到一块光轴平行于表面的方解石波片上，光的振动面和波片的主截面成 $30^\circ$ 角。求：（1）透射出来的寻常光和非常光的相对强度为多少？（2）用钠光入射时如要产生 $90^\circ$ 的相位差，波片的厚度至少应为多少？（$\lambda=589\,\mathrm{nm}$）}
\answerarea
\clearpage

\problem{5.9}{有一块平行石英片是沿平行于光轴方向切出的，要把它切成一块黄光的 $1/4$ 波片，问这块石英片应有多厚？石英的 $n_{\mathrm{e}}=1.552$，$n_{\mathrm{o}}=1.543$，$\lambda=589.3\,\mathrm{nm}$。}
\answerarea

\problem{5.10}{（1）线偏振光垂直入射到一个表面和光轴平行的波片，透射出来的光中，原来在波片中的寻常光及非常光产生了大小为 $\pi$ 的相位差，问波片厚度为多少？已知 $n_{\mathrm{o}}=1.544\,2$，$n_{\mathrm{e}}=1.553\,3$，$\lambda=500\,\mathrm{nm}$；（2）问这块波片应怎样放置才能使透射出来的光是线偏振光，而且它的振动面和入射光的振动面成 $90^\circ$ 角？}
\answerarea
\clearpage

\problem{5.11}{线偏振光垂直入射到一块表面平行于光轴的双折射晶片，光的振动面和晶片光轴成 $25^\circ$ 角，问波片中的寻常光和非常光透射出来后的相对强度如何？}
\answerarea

\problem{5.12}{在两个正交尼科耳棱镜 $\mathrm{N_1}$ 和 $\mathrm{N_2}$ 之间垂直插入一块波片，发现 $\mathrm{N_2}$ 后面有光射出。但当 $\mathrm{N_2}$ 绕入射光向顺时针转过 $20^\circ$ 后，$\mathrm{N_2}$ 的视场全暗。此时，把波片也绕入射光顺时针转过 $20^\circ$，$\mathrm{N_2}$ 的视场又亮了。问：（1）这是什么性质的波片；（2）$\mathrm{N_2}$ 要转过多大角度才能使 $\mathrm{N_2}$ 的视场又变为全暗。}
\answerarea
\clearpage

\problem{5.13}{一束圆偏振光，（1）垂直入射到 $1/4$ 波片上，求透射光的偏振状态；（2）垂直入射到 $1/8$ 波片上，求透射光的偏振状态。}
\answerarea

\problem{5.14}{试证明一束左旋圆偏振光和一束右旋圆偏振光，当它们的振幅相等时，合成的光是线偏振光。}
\answerarea
\clearpage

\problem{5.15}{设一方解石晶片沿平行光轴方向切出，其厚度为 $0.034\,3\,\mathrm{mm}$，放在两个正交的尼科耳棱镜间。平行光束经过第一尼科耳棱镜后，垂直地射到波片上，对于钠光（$589.3\,\mathrm{nm}$）而言，晶体的折射率为 $n_{\mathrm{o}}=1.658$，$n_{\mathrm{e}}=1.486$。试问：（1）通过第二尼科耳棱镜后，光束发生的干涉是加强还是减弱？（2）如果两个尼科耳棱镜的主截面是互相平行的，结果又如何？}
\answerarea

\problem{5.16}{单色光通过一尼科耳棱镜 $\mathrm{N_1}$，然后射到杨氏干涉实验装置的两个细缝上，试问：（1）尼科耳棱镜 $\mathrm{N_1}$ 的主截面与图面应成怎样的角度才能使光屏上的干涉图样中的暗条纹为最暗？（2）在上述情况下，在一个细缝前放置一半波片，并将这半波片绕着光线方向继续旋转，问在光屏上的干涉图样有何改变？}
\answerarea
\clearpage

\problem{5.17}{单色平行自然光垂直入射在杨氏双缝上，屏幕上出现一组干涉条纹。已知屏上 $A$、$C$ 两点分别对应零级亮纹和零级暗纹，$B$ 是 $AC$ 的中点，如题 5.17 图所示，试问：（1）若在双缝后放一理想偏振片 $\mathrm{P}$，屏上干涉条纹的位置、宽度会有何变化？$A$、$C$ 两点的光强会有何变化？（2）在一条缝的偏振片后放一片光轴与偏振片透光方向成 $45^\circ$ 的半波片，屏上有无干涉条纹？$A$、$B$、$C$ 各点的情况如何？\fig[0.52]{5_17}}
\answerarea

\problem{5.18}{厚度为 $0.025\,\mathrm{mm}$ 的方解石晶片，其表面平行于光轴，放在两个正交的尼科耳棱镜之间，光轴与两个尼科耳棱镜的主截面各成 $45^\circ$。如果射入第一个尼科耳棱镜的光是波长为 $400\sim 760\,\mathrm{nm}$ 的可见光，问透过第二个尼科耳棱镜的光中，少了哪些波长的光？}
\answerarea
\clearpage

\sproblem{5.19}{把一块切成长方体的 KDP 晶体放在两个正交的偏振片之间组成一个产生泡克耳斯效应的装置。已知电光常量 $\gamma=1.06\times 10^{-11}\,\mathrm{m/V}$，寻常光在该晶体中的折射率 $n_{\mathrm{o}}=1.51$，若入射光波长为 $550\,\mathrm{nm}$，试计算从晶体出射的两束线偏振光相位差为 $\pi$ 时，所需加在晶体上的纵向电压（叫做半波电压）。}
\answerarea

\sproblem{5.20}{将厚度为 $1\,\mathrm{mm}$ 且垂直于光轴切出的石英片放在两个平行的尼科耳棱镜之间，使从第一个尼科耳棱镜出射的光垂直射到石英片上，某一波长的光波经此石英片后，振动面旋转了 $20^\circ$。问石英片厚度为多少时，该波长的光将完全不能通过？}
\answerarea
\clearpage

\sproblem{5.21}{试求使波长为 $509\,\mathrm{nm}$ 的光的振动面旋转 $150^\circ$ 的石英片厚度，已知石英对这种波长的旋光度为 $29.7^\circ\,\mathrm{mm^{-1}}$。}
\answerarea

\sproblem{5.22}{将某种糖配制成浓度不同的 $4$ 种溶液：$100\,\mathrm{cm^3}$ 溶液中分别含有 $30.5\,\mathrm{g}$、$22.76\,\mathrm{g}$、$29.4\,\mathrm{g}$ 和 $17.53\,\mathrm{g}$ 的糖。分别用旋光量糖计测出它们通过每分米溶液转过的角度依次是 $49.5^\circ$、$36.1^\circ$、$30.3^\circ$ 和 $26.8^\circ$。根据这些结果，算出这几种糖的旋光率的平均值是多少？}
\answerarea
\clearpage

\sproblem{5.23}{如题 5.23 图所示装置中，$S$ 为单色光源，置于透镜 $\mathrm{L}$ 的焦点处，$\mathrm{P}$ 为起偏器，$\mathrm{L_1}$ 为此单色光的 $1/4$ 波片，其光轴与偏振器的透振方向成 $\alpha$ 角，$\mathrm{M}$ 为平面反射镜。已知入射到偏振器的光束强度为 $I_0$，试通过分析光束经过各元件后的光振动状态，求出光束返回后的光强 $I$。各元件对光束的损耗可忽略不计。\fig[0.52]{5_23}}
\answerarea

\sproblem{5.24}{一束椭圆偏振光沿 $z$ 方向传播，通过一个线起偏器。当起偏器透振方向沿 $x$ 方向时，透射强度最大，其值为 $1.5I_0$；当透振方向沿 $y$ 方向时，透射强度最小，其值为 $I_0$。（1）当透振方向与 $x$ 轴成 $\theta$ 角时，透射强度为多少？（2）使原来的光束先通过一个 $1/4$ 波片后再通过线起偏器，$1/4$ 波片的轴沿 $x$ 方向。现在发现，当起偏器透光方向与 $x$ 轴成 $30^\circ$ 角时，透过两个元件的光强最大，求光强的最大值，并确定入射光强中非偏振成分占多少？}
\answerarea
\clearpage

\sproblem{5.25}{有下列几个未加标明的光学元件：（1）两个线偏振器；（2）一个 $1/4$ 波片；（3）一个半波片；（4）一个圆偏振器。除了一个光源和一个光屏外不借助其他光学仪器，如何鉴别上述各个元件。}
\answerarea

\sproblem{5.26}{一束汞绿光以 $60^\circ$ 角入射到磷酸二氢钾（KDP）晶体表面，晶体的 $n_{\mathrm{o}}=1.512$，$n_{\mathrm{e}}=1.470$。设光轴与晶面平行，并垂直于入射面，试求晶体中 $\mathrm{o}$ 光与 $\mathrm{e}$ 光的夹角。}
\answerarea
\clearpage

\sproblem{5.27}{通过尼科耳棱镜观察一束椭圆偏振光时，强度随尼科耳棱镜的旋转而改变。当强度为极小值时，在尼科耳棱镜（检偏器）前插入一块 $1/4$ 波片，转动 $1/4$ 波片使它的光轴平行于检偏器的透振方向，再把检偏器沿顺时针方向转动 $20^\circ$ 就完全消光。问（1）该椭圆偏振光是右旋还是左旋的？（2）椭圆长、短轴之比是多少？}
\answerarea

\sproblem{5.28}{研究性课题：请查阅互联网，了解偏振式 3D 技术的新进展，并写成一篇小论文。}
\answerarea
\clearpage

\chapter{光的吸收、散射和色散}

\problem{6.1}{一固体有两个吸收带，宽度都是 $30\,\mathrm{nm}$。一带处在蓝光区（$450\,\mathrm{nm}$ 附近），另一带处在黄光区（$580\,\mathrm{nm}$ 附近）。设第一带吸收系数为 $50\,\mathrm{cm^{-1}}$，第二带的为 $250\,\mathrm{cm^{-1}}$。试绘出白光分别透过 $0.1\,\mathrm{mm}$ 及 $5\,\mathrm{mm}$ 的该物质后在吸收带附近光强分布的概况。}
\answerarea

\problem{6.2}{某种介质的吸收系数 $\alpha_{\mathrm{a}}$ 为 $0.32\,\mathrm{cm^{-1}}$，求透射光强为入射光强的 $0.1$、$0.2$、$0.5$ 及 $0.8$ 倍时，该介质的厚度各为多少？}
\answerarea
\clearpage

\problem{6.3}{如果同时考虑到吸收和散射都将使透射光强度减弱，则透射光表达式中的 $\alpha$ 可看做是由两部分合成，一部分 $\alpha_{\mathrm{a}}$ 是由于真正的吸收（变为物质分子的热运动），另一部分 $\alpha_{\mathrm{s}}$（称为散射系数）是由于散射，于是该式可写作 $I=I_0\mathrm{e}^{-(\alpha_{\mathrm{a}}+\alpha_{\mathrm{s}})d}$。如果光通过一定厚度的某种物质后，只有 $20\%$ 的光强通过，已知该物质的散射系数等于吸收系数的 $1/2$。假定不考虑散射，则透射光强可增加多少？}
\answerarea

\problem{6.4}{计算波长为 $253.6\,\mathrm{nm}$ 和 $546.1\,\mathrm{nm}$ 的两条谱线瑞利散射的强度之比。}
\answerarea
\clearpage

\problem{6.5}{太阳光束由小孔射入暗室，室内的人沿着与光束垂直及与之成 $45^\circ$ 的方向观察这光束时，见到瑞利散射的散射光强度之比等于多少？}
\answerarea

\problem{6.6}{一束光通过液体，用尼科耳棱镜正对这束光进行观察，当尼科耳棱镜主截面竖直时，光强达最大值；当尼科耳棱镜主截面水平时，光强为零。再从侧面观察其散射光，在尼科耳棱镜主截面为竖直和水平两个位置时，光强之比为 $20:1$，计算散射光的退偏振度。}
\answerarea
\clearpage

\problem{6.7}{一种光学玻璃对于波长 $435.8\,\mathrm{nm}$ 和 $546.1\,\mathrm{nm}$ 的折射率分别为 $1.613\,0$ 和 $1.602\,6$，试应用柯西公式来计算这种玻璃对波长 $600\,\mathrm{nm}$ 的光的色散 $\dfrac{\mathrm{d}n}{\mathrm{d}\lambda}$。}
\answerarea

\problem{6.8}{一种光学玻璃对汞蓝光 $435.8\,\mathrm{nm}$ 和汞绿光 $546.1\,\mathrm{nm}$ 的折射率分别为 $1.652\,50$ 和 $1.624\,50$。（1）试用柯西公式计算公式中的常量 $a$ 和 $b$；（2）并求它对 $589\,\mathrm{nm}$ 钠黄光的折射率 $n$ 和色散 $\dfrac{\mathrm{d}n}{\mathrm{d}\lambda}$。}
\answerarea
\clearpage

\problem{6.9}{一个顶角为 $60^\circ$ 的棱镜由某种玻璃制成，它的色散特性可用柯西公式中的常量 $a=1.416$，$b=1.72\times 10^{-10}\,\mathrm{cm^2}$ 来表示。将棱镜的位置放置得使它对 $600\,\mathrm{nm}$ 的波长产生最小偏向角，试计算这个棱镜的角色散率为多少。}
\answerarea

\problem{6.10}{波长为 $0.67\,\mathrm{nm}$ 的 X 射线，由真空入射到某种玻璃时，在掠射角不超过 $0.1^\circ$ 的条件下发生全反射，试计算玻璃对这波长的折射率，并解释所得的结果。}
\answerarea
\clearpage

\sproblem{6.11}{波长为 $0.07\,\mathrm{nm}$ 的 X 射线，其折射率比 $1$ 小 $1.60\times 10^{-6}$，试求全反射时，X 射线的掠射角为多大？}
\answerarea

\sproblem{6.12}{对于波长 $\lambda=400\,\mathrm{nm}$ 的光，某种玻璃的折射率 $n=1.63$，对 $\lambda'=500\,\mathrm{nm}$ 的光，其折射率 $n'=1.58$。假如柯西公式的近似形式 $n=a+\dfrac{b}{\lambda^2}$ 适用，试求此种玻璃对波长为 $600\,\mathrm{nm}$ 的光的色散 $\dfrac{\mathrm{d}n}{\mathrm{d}\lambda}$ 值。}
\answerarea
\clearpage

\sproblem{6.13}{研究性课题：$\mathrm{PM2.5}$ 与大气环境污染检测。}
\answerarea
\clearpage

\chapter{光的量子性}

\problem{7.1}{在深度远大于表面波波长的液体中，表面波的传播速度满足如下规律：
\[v=\sqrt{\frac{\lambda}{2\pi}\left(g+\frac{4\pi^2 F}{\lambda^2\rho}\right)}\]
式中 $g$ 为重力加速度，$\rho$ 为液体密度，$F$ 为表面张力，$\lambda$ 为表面波的波长。试计算表面波的群速度。}
\answerarea

\problem{7.2}{测量二硫化碳的折射率的实验数据为：当 $\lambda'=589\,\mathrm{nm}$，$n'=1.629$；当 $\lambda''=656\,\mathrm{nm}$ 时，$n''=1.620$。试求波长 $589\,\mathrm{nm}$ 的光在二硫化碳中的相速、群速和群速折射率。}
\answerarea
\clearpage

\problem{7.3}{在测定光速的迈克耳孙旋转棱镜法中，设所用棱镜为正 $n$ 面棱柱体。试导出：根据棱镜的转速、反射镜距离等数据计算光速的公式。}
\answerarea

\problem{7.4}{试用光的相速度 $v$ 和 $\dfrac{\mathrm{d}v}{\mathrm{d}\lambda}$ 来表示群速度 $u=\dfrac{\mathrm{d}\omega}{\mathrm{d}k}$，再用 $v$ 和 $\dfrac{\mathrm{d}n}{\mathrm{d}\lambda}$ 表示群速度 $u=\dfrac{\mathrm{d}\omega}{\mathrm{d}k}$。}
\answerarea
\clearpage

\problem{7.5}{计算在下列各种色散介质中传播的各种不同性质的波的群速度：（1）$v=\text{常量}$（无色散介质，如空气中的声波）；（2）$v=a\sqrt{\lambda}$，$a$ 为常量（重力在水面上所引起的波）；（3）$v=\dfrac{a}{\sqrt{\lambda}}$（在水面上的表面张力波）；（4）$v=\dfrac{a}{\lambda}$（弹性薄片在弯曲时所产生的波）；（5）$v=\sqrt{c^2+b^2\lambda^2}$（电离层中的电磁波，其中 $c$ 是真空中的光速，$\lambda$ 是介质中的波长）；（6）$v=c\omega/\sqrt{\omega^2\varepsilon\mu-c^2a^2}$（在充满色散介质的直波导管中的电磁波，式中 $c$ 是真空中的光速，$a$ 是与波导管横截面的形状及大小有关的常量，$\varepsilon=\varepsilon(\omega)$ 是介质的电容率，$\mu=\mu(\omega)$ 是介质的磁导率）。}
\answerarea

\problem{7.6}{利用维恩公式求：辐射的最概然频率 $\nu_{\mathrm{m}}$、辐射的最概然波长 $\lambda_{\mathrm{m}}$、辐射的最大光谱密度 $(\varepsilon_\lambda)_{\mathrm{m}}$、辐射出射度 $M_0(T)$ 与温度 $T$ 的关系。}
\answerarea
\clearpage

\problem{7.7}{太阳光谱非常接近于 $\lambda_{\mathrm{m}}=480\,\mathrm{nm}$ 的绝对黑体的光谱。试求：（1）在 $1\,\mathrm{s}$ 内太阳由于辐射而损失的质量；（2）并估计太阳的质量减少 $1\%$（由于热辐射）所经历的时间。太阳的质量 $m_0$ 为 $2.0\times 10^{30}\,\mathrm{kg}$，太阳的半径 $r$ 是 $7.0\times 10^8\,\mathrm{m}$。}
\answerarea

\problem{7.8}{地球表面每平方厘米每分钟由于辐射而损失的能量平均值为 $0.546\,\mathrm{J}$。如有一个黑体，它在辐射相同的能量时，温度应为多少？}
\answerarea
\clearpage

\problem{7.9}{若有一黑体的辐出度等于 $5.70\,\mathrm{W/cm^2}$，试求与该辐射最大光谱强度相对应的波长。}
\answerarea

\problem{7.10}{用交流供电的灯丝温度是变动的。一电灯白炽时钨丝的平均温度为 $2\,300\,\mathrm{K}$，其最高和最低温度之差约为 $80\,\mathrm{K}$。问热辐射的总功率的最大值和最小值之比为多少？钨丝的辐射可当做黑体。}
\answerarea
\clearpage

\problem{7.11}{若将恒星表面的辐射近似地看做是黑体辐射，则可以用测量 $\lambda_{\max}$ 的办法来估计恒星表面的温度。现测得太阳的 $\lambda_{\max}$ 为 $510\,\mathrm{nm}$，北极星的 $\lambda_{\max}$ 为 $350\,\mathrm{nm}$，试求它们的表面温度。}
\answerarea

\problem{7.12}{小灯泡所耗的功率为 $1\,\mathrm{W}$，均匀地向各个方向辐射能量。设辐射的平均波长为 $500\,\mathrm{nm}$，试求在 $10\,\mathrm{km}$ 处每秒钟落在垂直于光线方向上每平方厘米面积上的光子数。}
\answerarea
\clearpage

\problem{7.13}{已知铯的逸出功为 $1.88\,\mathrm{eV}$，现用波长为 $300\,\mathrm{nm}$ 的紫外光照射，试求光电子的初动能和初速。}
\answerarea

\problem{7.14}{用波长为 $253\,\mathrm{nm}$ 的光照射钨的表面时，在光电管的电路中产生的光电流，由于外加 $1\,\mathrm{V}$ 的遏止电压而截止。已知钨的逸出功为 $4.5\,\mathrm{eV}$，试确定接触电势差。}
\answerarea
\clearpage

\problem{7.15}{波长为 $320\,\mathrm{nm}$ 的紫外线入射到逸出功为 $2.2\,\mathrm{eV}$ 的金属表面上，求光电子从金属表面逸出时的最大速度。若入射光的波长为原来的一半，出射光电子的最大动能是否增至两倍？}
\answerarea

\problem{7.16}{波长为 $0.1\,\mathrm{nm}$ 的 X 射线被碳块散射，在散射角为 $90^\circ$ 的方向上进行观察。试求（1）康普顿位移 $\Delta\lambda$；（2）反冲电子的动能。}
\answerarea
\clearpage

\problem{7.17}{已知入射光子的波长为 $0.003\,\mathrm{nm}$，而反冲电子的速度为光速的 $\beta$ 倍（$\beta=0.6$），试确定康普顿位移 $\Delta\lambda$。}
\answerarea

\problem{7.18}{在电子显微镜中，电子受到 $90\,\mathrm{kV}$ 的电压加速，如要观察到物质的分子结构（其大小为 $10^{-9}\,\mathrm{cm}$ 数量级），问显微镜的数值孔径应该多大？}
\answerarea
\clearpage

\sproblem{7.19}{（1）一只 $100\,\mathrm{W}$ 灯泡，$5\%$ 的功率辐射是可见光，假定可见光平均波长为 $500\,\mathrm{nm}$，则每秒可辐射的可见光光子的能量为多少？（2）假设灯泡为点光源，可以向空间各个方向发光，求在距离 $2\,\mathrm{m}$ 处每秒垂直通过单位面积的光子数。}
\answerarea

\sproblem{7.20}{人在黑暗中，眼睛的视网膜如能接收到波长为 $550\,\mathrm{nm}$ 的最大有效辐射为 $2\times 10^{-18}\,\mathrm{J}$，就能感知这一点光源。求眼睛这一观察阈值相当于多少个光子？}
\answerarea
\clearpage

\sproblem{7.21}{一电子束被电压为 $V$ 的电场加速。（1）求电子在被加速后的德布罗意波长；（2）求此德布罗意波的相速和群速；（3）把此德布罗意波射到一块单晶上，如入射方向与晶面成 $\theta$ 角，观察到散射波第一级强度极大值，求晶格常量 $d$。}
\answerarea

\sproblem{7.22}{在热核爆炸的火球中，测得瞬时温度为 $10^7\,\mathrm{K}$。（1）估算辐射最强的波长；（2）这种辐射最强波长的能量子是多少焦耳？}
\answerarea
\clearpage

\sproblem{7.23}{用关键词光镊（optical tweezer）或激光捕获（laser trapping），在互联网（如 Google 或 Yahoo）上查阅光镊技术在生物细胞检测和医学诊断方面的应用和进展，写一篇介绍报告。}
\answerarea
\clearpage

\chapter{现代光学基础}

\problem{8.1}{（1）计算氢原子最低的四个能级的能量大小，并把它们画成能级图；（2）计算这四个能级之间跃迁的最小的频率是多少？}
\answerarea

\problem{8.2}{当玻尔描述的氢原子从 $n=2$ 的轨道跃迁到 $n=1$ 的轨道后，问（1）轨道半径变化了多少？（2）能量改变了多少？}
\answerarea
\clearpage

\problem{8.3}{一个电子受激跃迁到氢原子的第三玻尔轨道，试计算回到基态时有可能发射的光子的波长。}
\answerarea

\problem{8.4}{氦氖激光的单色性为 $6\times 10^{-10}$，则其相干长度为多少？}
\answerarea
\clearpage

\problem{8.5}{证明当每个模内的平均光子数大于 $1$ 时，腔内振荡以受激辐射为主。（提示：以受激辐射为主时，$\dfrac{w_{21}}{A_{21}}=\dfrac{B_{21}\rho}{A_{21}}>1$。）}
\answerarea

\problem{8.6}{一质地均匀的材料对光的吸收系数为 $0.01\,\mathrm{mm^{-1}}$，光通过 $10\,\mathrm{cm}$ 长的该材料后，光强损耗多少？若对 $10\,\mathrm{cm}$ 长的另一种材料，出射光强是入射光强的 $2.718$ 倍，问该材料的增益系数为多少？}
\answerarea
\clearpage

\problem{8.7}{如激光器分别在波长为 $10\,\mathrm{\mu m}$ 和 $500\,\mathrm{nm}$ 时输出 $1\,\mathrm{W}$ 的连续功率，求每秒从激光上能级向下能级跃迁时对应的粒子数。}
\answerarea

\problem{8.8}{设一对激光能级为 $E_2$ 和 $E_1$，上下能级上的粒子密度分别为 $n_2$ 和 $n_1$，两能级间的跃迁频率为 $\nu$，波长为 $\lambda$。试求（1）当 $\nu=3\,000\,\mathrm{MHz}$ 和 $\lambda=1\,\mathrm{\mu m}$ 时，$n_2/n_1$ 分别为多少？假定此时 $T=300\,\mathrm{K}$；（2）当 $\lambda=1\,\mathrm{\mu m}$，$n_2/n_1=0.1$ 时，温度 $T$ 为多少？}
\answerarea
\clearpage

\problem{8.9}{假若受激辐射爱因斯坦系数 $B=10^{19}\,\mathrm{m^3/(W\cdot s^3)}$，试对下列波长计算自发辐射跃迁系数和自发辐射寿命。（1）$\lambda=6\,\mathrm{\mu m}$；（2）$\lambda=600\,\mathrm{nm}$；（3）$\lambda=60\,\mathrm{nm}$；（4）$\lambda=0.6\,\mathrm{nm}$。如果光强为 $10\,\mathrm{W/mm^2}$，求受激跃迁速率。}
\answerarea

\problem{8.10}{有一个二能级系统，如果能级差为 $1.602\times 10^{-21}\,\mathrm{J}$，求：（1）当 $T=10^2\,\mathrm{K}$、$10^5\,\mathrm{K}$ 和 $10^8\,\mathrm{K}$ 时，上下两能级粒子数 $N_2$ 与 $N_1$ 之比；（2）如 $N_2=N_1$，则相当于多高的温度 $\mathrm{K}$；（3）要有怎样的温度，粒子数才能发生反转；（4）如果用“负温度”的概念来描述粒子数反转的状态，则 $T=-10^4\,\mathrm{K}$ 和 $T=-10^8\,\mathrm{K}$ 两个温度中哪一个高？$T=-10^5\,\mathrm{K}$ 和 $T=+10^8\,\mathrm{K}$ 中哪一个高？}
\answerarea
\clearpage

\problem{8.11}{衍射和傅里叶变换的关系是什么？如何应用傅里叶变换来计算菲涅耳衍射和夫琅禾费衍射？给出计算的方法和步骤。}
\answerarea

\problem{8.12}{从惠更斯-菲涅耳原理的数学表达式出发，
\begin{enumerate}
  \item 求出菲涅耳衍射方程；
  \item 证明菲涅耳衍射的作用相当于空间不变的线性系统；
  \item 并求出传递函数；
  \item 求出夫琅禾费衍射方程。
\end{enumerate}}
\answerarea
\clearpage

\problem{8.13}{描述二元振幅滤波器的种类及其应用。}
\answerarea

\problem{8.14}{说明泽尼克相衬法中采用的滤波系统，并计算像面上的强度分布情况。}
\answerarea
\clearpage

\problem{8.15}{题 8.15 图为电寻址液晶光阀试验系统，主要由激光光源、准直扩束系统、高分辨率电寻址透射式液晶光阀、傅里叶变换透镜、CCD 探测器等构成。与光寻址液晶光阀不同的是，该系统可以直接把计算机产生的信息显示在液晶光阀上，可以实时地进行图像处理，且方便实验操作。

请解释该系统的工作原理。\figcap[0.74]{8_15}{题 8.15 图\quad 电寻址液晶光阀实验系统}}
\answerarea

\sproblem{8.16}{研究性课题：通过书面材料或互联网，查找激光技术的进展，或激光的应用等方面的知识，写出报告，在同学之间相互交流。}
\answerarea
\clearpage

\end{document}
